%% Extraído de docs/TCC_documento/Textuais/Metodologia/08_SAME.tex (linhas 239-252)
\begin{tikzpicture}[>=Latex]
    \draw[fill=black!12, draw=black, thick]
        (100:2.0) arc (100:440:2.0) -- (80:1.3) arc (80:-260:1.3) -- cycle;
    \draw[->, thick, blue!55!black] (200:1.65) arc (200:260:1.65);
    \draw[->, thick, blue!55!black] (280:1.65) arc (280:340:1.65);
    \node[blue!55!black] at (0,-1.62) {$\Phi$};
    \draw[thick, ->] (-2.6,0) -- (4.7,0) node[right] {$i(t)$};
    \node[above, font=\footnotesize] at (3.45,0.02) {condutor de potência};
    \fill[black] (-0.15,1.48) rectangle (0.15,1.82);
    \draw[gray] (0.15,1.68) -- (1.1,2.15);
    \node[right] at (1.1,2.15) {Sensor Hall};
    \draw[gray] (145:2.0) -- (-2.5,1.35);
    \node[left] at (-2.5,1.35) {Núcleo toroidal (ferrite)};
\end{tikzpicture}
